% !TeX root = ../main.tex
窗外的桂花又开了，淡淡的花香随着秋风飘入室内，才恍然意识到，秋天又一次来到了身边。时光匆匆，今年已是我来到西安交通大学的第六年。六年前，我从南京乘坐一夜火车来到西安，从一座古都来到另一座古都。六年间，我逐渐熟悉并喜欢上这座城市的饮食、气候与生活节奏，也在这里度过了人生中一段充实而难忘的时光。如今，求学生活即将告一段落，回望一路走来的点滴，心中更多的是感激。

首先，衷心感谢我的导师恽迪教授。恽老师在科研和学习中给予我充分的信任与自主空间，尊重学生的成长节奏，也始终在关键问题上给予耐心指导。无论是科研方向的把握、具体问题的分析，还是论文写作与修改，老师的指导都使我受益良多。在这样的学习氛围中，我逐渐养成了规律学习、主动思考和自我管理的习惯，也学会在科研之外通过运动、电影等方式调节身心。老师严谨而包容的治学态度，对我今后的学习和工作都会产生长久影响。

感谢柳文波老师对论文、组会汇报和周报的细致指导，使我能够及时总结每周的收获与不足，在不断反思中改进科研工作。感谢雷鹏辉老师在求职过程中的关心和帮助。由衷感谢单建强、苟军利、邱杰、张斌、吴攀、葛莉、{\rarecharfont 桂淼}等老师在学习、科研和生活中给予的支持与帮助。感谢课题组历任组长组织丰富的春游、晚会、温泉和城墙活动，这些共同经历缓解了科研生活的紧张，也为我的求学生涯留下了许多温暖的回忆。

感谢王兆、张新虎师兄在科研入门阶段给予的帮助；感谢钱郑宇、王挺、王冠、袁和师兄在我遇到学术问题时给予耐心细致的解答；感谢张晋恺师弟在MOOSE学习和使用过程中的相互交流与帮助；感谢冯琳娜、师田田在职业规划方面给予的建议；感谢葛震鸿师弟在论文修改的建议。感谢王扬中、文春阳、{\rarecharfont 郝喆}、崔哲鑫、姜彦博、严冬、徐钰雯、{\rarecharfont 赵祎玮}、戚晓勇、兰峋、佘仲航、刘续希、喇永孝、郑诗伦、孟睿智、赵芬、高士森、吴佳轩、杨重斗、雷敏、艾靖凯、陈可心、郑睿、翟静杰、李思雯等同学在学习、科研和生活中给予的帮助与陪伴。

感谢我的父母。虽然家庭条件并不宽裕，但他们始终尽己所能支持我的学习和选择，给予我最坚定、最无私的理解与鼓励。正是这份支持，使我能够安心完成学业并不断前行。

谨向所有曾经关心、帮助和陪伴过我的老师、同学、朋友和家人致以诚挚的谢意。
